\begin{table}[tp]
\centering
\caption{Oracle tradeoffs for \href{https://github.com/k3nfalt/longtails/blob/main/LongTails/Setups/ThirdDerivativeSmoothNonconvex.lean}{Third\allowbreak{}Derivative\allowbreak{}Smooth\allowbreak{}Nonconvex}.}\label{tab:ThirdDerivativeSmoothNonconvex}
\parbox{\linewidth}{\footnotesize On a finite-dimensional real inner-product space, $f\in C^3$ has $L_3$-Lipschitz third derivative, $f(x_0)-f(x)\le\Delta$ for every $x$, and $\|\nabla^2 f(x_0)\|\le H_0$, with $L_3,\Delta,\varepsilon>0$ and $H_0\ge0$. Exact $(f,\nabla f)$ and full-Hessian oracles are available. The target is $\|\nabla f(x_{\mathrm{out}})\|\le\varepsilon$; \textbf{(1)} guarantees it for every seed.}\par\smallskip
\begin{threeparttable}
\footnotesize
\setlength{\tabcolsep}{4pt}\renewcommand{\arraystretch}{1.25}
\begin{tabularx}{\linewidth}{@{}>{\raggedright\arraybackslash}p{0.20\linewidth}>{\raggedright\arraybackslash}p{0.32\linewidth}>{\raggedright\arraybackslash}p{0.25\linewidth}>{\raggedright\arraybackslash}X@{}}
\toprule
\textbf{Result} & \textbf{Oracle complexity} & \textbf{Novelty} & \textbf{Limitations} \\
\midrule
\textbf{(1)} \href{https://github.com/k3nfalt/longtails/blob/main/LongTails/Theorems/LazyQuarticHessianTradeoff.lean}{Lazy\allowbreak{}Quartic\allowbreak{}Hessian\allowbreak{}Tradeoff} \textbf{(Ours)} $(\star\star\star)$ \href{https://github.com/k3nfalt/longtails/blob/main/scripts/reviews/LazyQuarticHessianTradeoff.md}{(review)}
& $\mathcal{O}\!\left((1+S^{5/4})\log\Gamma\right.$\newline
$\left.{}+m(1+S^{3/4})\right)$
& --- & --- \\
\midrule[1.2pt]
Lower-order tensor method \citep[Thm.~3.5, Cor.~3.6, $p=3$]{tdsnDoikov2025v1}
& $\mathcal{O}\!\left(m\left[d^{1/3}\Delta\widehat L_3^{1/3}\varepsilon^{-4/3}\right.\right.$\newline
$\left.\left.{}+d\bigl(1+\log_+(L_3/L_{\mathrm{init}})\bigr)\right]\right)$\newline
Value/gradient and Hessian calls, including trials.
& \textbf{(1)}: i) Dimension-free; Hessians $S^{3/4}$ versus $d^{1/3}S+d$ for $S\ge1$.\newline
 ii) Total gain when $L_{\mathrm{init}}\asymp L_3$, $S\ge1$, and $md^{1/3}\gg S^{1/4}\log\Gamma$.
& \textbf{(1)}: i) More gradients when $S^{1/4}\log\Gamma\gg d^{1/3}$; worse for $m=d=1$, $S\gg1$.\newline
 ii) Needs $L_3,\Delta,H_0$; baseline adapts to unknown $L_3$. \\
\midrule
Cutting planes \citep[Thm.~1, first-order variant after (8)]{tdsnHinder2018v2}\newline
Special case: globally $L_1$-Lipschitz gradient.
& $\mathcal{O}\!\left(d(1+S)\left[1+\log_+\!\frac{L_1^3}{L_3\varepsilon^2}\right]\right)$\newline
Value/gradient calls; no Hessians.
& \textbf{(1)}: i) Scope: no global $L_1$; dimension-free.\newline
 ii) On overlap, total gain for $S\ge1$, $d\gg S^{1/4}$, $m\ll dS^{1/4}$, up to logarithms.
& \textbf{(1)}: Needs Hessians; worse polynomial bound when $d\lesssim S^{1/4}$, $S\gg1$. \\
\midrule
Guarded AGD \citep[Thm.~2]{tdsnCarmon2017}\newline
Special case: globally $L_1$-Lipschitz gradient.
& $\mathcal{O}\!\left(\Delta L_1^{1/2}L_3^{1/6}\varepsilon^{-5/3}\right.$\newline
$\left.{}\cdot[1+\log_+(L_1\Delta/\varepsilon^2)]\right)$\newline
Value/gradient calls, provided\newline
$\varepsilon^{2/3}\le\delta^{1/2}L_3^{1/6}$,\newline
$\varepsilon^{2/3}\le L_1/(8L_3^{1/3})$.
& \textbf{(1)}: i) Scope: no global $L_1$.\newline
 ii) On overlap, replaces $\sqrt{L_1}$ by $(L_3\Delta)^{1/4}$; total gain for $L_1\gg\sqrt{L_3\Delta}$, $m\ll\Delta^{1/4}L_1^{1/2}L_3^{-1/12}\varepsilon^{-2/3}$, up to logarithms.
& \textbf{(1)}: Needs Hessians and $H_0$ in addition; worse polynomial bound when $L_1\lesssim\sqrt{L_3\Delta}$. \\
\bottomrule
\end{tabularx}
\begin{tablenotes}[flushleft]\footnotesize
\item[] \textbf{Abbr:} $G,H$: numbers of value/gradient and full-Hessian calls; displayed costs are $G+mH$, with Hessian cost $m\ge1$ in value/gradient units. $d=\dim E$; $S=\Delta L_3^{1/3}\varepsilon^{-4/3}$; $\Gamma=2+H_0\Delta/\varepsilon^2+L_3\Delta^3/\varepsilon^4$; $\delta=f(x_0)-\inf f$. $L_1$ is a global gradient Lipschitz constant; $L_{\mathrm{init}}>0$ is the tensor method's initial smoothness estimate; $\widehat L_3=\max\{L_{\mathrm{init}},L_3\}$. $\log_+t=\max\{0,\log t\}$; $\asymp$ denotes equality up to universal factors.
\end{tablenotes}
\end{threeparttable}
\end{table}
